% ============================================================
\chapter{Espaces de Sobolev}
\label{ch:sobolev}
% ============================================================

Les \'equations aux d\'eriv\'ees partielles classiques exigent que les solutions soient suffisamment r\'eguli\`eres pour que les d\'eriv\'ees aient un sens ponctuel. Mais de nombreux probl\`emes physiques --- membranes vibrantes avec des coins, \'ecoulements turbulents, mat\'eriaux composites --- produisent des solutions qui ne sont pas diff\'erentiables au sens classique. Sergei Sobolev, dans les ann\'ees 1930, a r\'esolu ce dilemme en introduisant les \emph{d\'eriv\'ees faibles}~: une notion de d\'erivation qui ne demande pas la r\'egularit\'e ponctuelle mais seulement que la \og formule d'int\'egration par parties \fg{} soit satisfaite au sens des distributions. Les \emph{espaces de Sobolev} $W^{k,p}(\Omega)$, munis de normes naturelles, sont les espaces fonctionnels dans lesquels vivent les solutions faibles des EDP. Aujourd'hui, toute la th\'eorie moderne des EDP --- m\'ethodes variationnelles, \'el\'ements finis, r\'egularit\'e --- repose sur ce cadre.

\section{D\'eriv\'ees faibles}

\begin{definition}[D\'eriv\'ee faible]
\index{d\'eriv\'ee faible}
Soit $\Omega \subset \R^n$ un ouvert et $u \in L^1_{\mathrm{loc}}(\Omega)$.
On dit que $v \in L^1_{\mathrm{loc}}(\Omega)$ est la \emph{d\'eriv\'ee faible}
$\partial^\alpha u$ (pour un multi-indice $\alpha$) si~:
\[
    \int_\Omega v\, \varphi\, dx = (-1)^{\abs{\alpha}} \int_\Omega u\, \partial^\alpha \varphi\, dx
    \quad \forall\, \varphi \in C_c^\infty(\Omega).
\]
\end{definition}

\begin{intuition}
La d\'eriv\'ee faible g\'en\'eralise la d\'eriv\'ee classique aux fonctions qui ne sont
pas n\'ecessairement diff\'erentiables au sens usuel. La formule repose sur une
int\'egration par parties~: si $u$ est $C^k$, la d\'eriv\'ee faible co\"incide avec
la d\'eriv\'ee classique. L'int\'er\^et est que la d\'eriv\'ee faible peut exister m\^eme
si $u$ a des singularit\'es.
\end{intuition}

\begin{exemple}[Valeur absolue]
La fonction $u(x) = \abs{x}$ sur $\R$ a pour d\'eriv\'ee faible
$u'(x) = \operatorname{sgn}(x)$ (fonction signe). En effet, pour tout
$\varphi \in C_c^\infty(\R)$~:
\[
    \int_\R \abs{x} \varphi'(x)\, dx = -\int_\R \operatorname{sgn}(x)\, \varphi(x)\, dx.
\]
\end{exemple}

\begin{exemple}[Fonction de Heaviside]
La fonction de Heaviside $H(x) = \ind_{(0,\infty)}(x)$ n'a pas de d\'eriv\'ee faible
dans $L^1_{\mathrm{loc}}(\R)$~: sa d\'eriv\'ee au sens des distributions est la
masse de Dirac $\delta_0$, qui n'est pas une fonction.
\end{exemple}

\section{Espaces de Sobolev $W^{k,p}$}

\begin{definition}[Espace de Sobolev $W^{k,p}(\Omega)$]
\index{espace de Sobolev}
Soit $\Omega \subset \R^n$ un ouvert, $k \in \N$ et $1 \leq p \leq \infty$.
L'\emph{espace de Sobolev} $W^{k,p}(\Omega)$ est~:
\[
    W^{k,p}(\Omega) = \{ u \in L^p(\Omega) : \partial^\alpha u \in L^p(\Omega),\;
    \forall\, \abs{\alpha} \leq k \},
\]
muni de la norme~:
\[
    \norm{u}_{W^{k,p}} = \left( \sum_{\abs{\alpha} \leq k}
    \norm{\partial^\alpha u}_{L^p}^p \right)^{1/p}
    \quad (1 \leq p < \infty).
\]
Pour $p = 2$, on note $H^k(\Omega) = W^{k,2}(\Omega)$.
\end{definition}

\begin{theoreme}[Compl\'etude]
L'espace $W^{k,p}(\Omega)$ est un espace de Banach. Pour $p = 2$, $H^k(\Omega)$
est un espace de Hilbert avec le produit scalaire~:
\[
    \inner{u}{v}_{H^k} = \sum_{\abs{\alpha} \leq k} \int_\Omega
    \partial^\alpha u \cdot \overline{\partial^\alpha v}\, dx.
\]
\end{theoreme}

\begin{definition}[Espace $W_0^{k,p}(\Omega)$]
\index{espace $W_0^{k,p}$}
L'espace $W_0^{k,p}(\Omega)$ est la fermeture de $C_c^\infty(\Omega)$ dans
$W^{k,p}(\Omega)$. Les fonctions de $W_0^{k,p}$ sont celles qui \og{}s'annulent
au bord\fg au sens de la trace. On note $H_0^k(\Omega) = W_0^{k,2}(\Omega)$.
\end{definition}

\section{Densit\'e des fonctions lisses}

\begin{theoreme}[Meyers-Serrin]
\index{th\'eor\`eme de Meyers-Serrin}
L'espace $C^\infty(\Omega) \cap W^{k,p}(\Omega)$ est dense dans $W^{k,p}(\Omega)$
pour $1 \leq p < \infty$.
\end{theoreme}

\begin{proof}[Id\'ee de preuve]
On utilise une partition de l'unit\'e $\{\rho_j\}$ et des r\'egularisations par
convolution $(\rho_j u) * \eta_{\varepsilon_j}$ avec des noyaux mollificateurs
$\eta_\varepsilon$. En choisissant $\varepsilon_j$ assez petit pour chaque $j$,
on approche $u$ dans $W^{k,p}$.
\end{proof}

\begin{theoreme}[Densit\'e de $C^\infty(\overline{\Omega})$]
Si $\Omega$ est un ouvert born\'e \`a bord $C^1$, alors $C^\infty(\overline{\Omega})$
est dense dans $W^{k,p}(\Omega)$ pour $1 \leq p < \infty$.
\end{theoreme}

\section{Th\'eor\`emes d'injection de Sobolev}

\begin{theoreme}[Injection de Sobolev]
\label{thm:sobolev_injection}
\index{injection de Sobolev}
Soit $\Omega \subset \R^n$ un ouvert born\'e \`a bord lipschitzien, $k \in \N^*$
et $1 \leq p < \infty$.
\begin{enumerate}
    \item Si $kp < n$~: $W^{k,p}(\Omega) \hookrightarrow L^{p^*}(\Omega)$
          avec $p^* = \frac{np}{n - kp}$ (exposant de Sobolev).
    \item Si $kp = n$~: $W^{k,p}(\Omega) \hookrightarrow L^q(\Omega)$
          pour tout $q \in [p, \infty)$.
    \item Si $kp > n$~: $W^{k,p}(\Omega) \hookrightarrow C^{0,\gamma}(\overline{\Omega})$
          avec $\gamma = k - n/p$ si $k - n/p < 1$.
\end{enumerate}
\end{theoreme}

\begin{attention}[title=\textbf{L'exposant critique}]
L'injection $W^{1,p} \hookrightarrow L^{p^*}$ est \emph{continue} mais
\emph{non compacte}. C'est le th\'eor\`eme de Rellich-Kondrachov qui donne
la compacit\'e pour les exposants strictement inf\'erieurs \`a $p^*$.
\end{attention}

\begin{theoreme}[Rellich-Kondrachov]
\label{thm:rellich}
\index{th\'eor\`eme de Rellich-Kondrachov}
Soit $\Omega$ un ouvert born\'e \`a bord lipschitzien.
\begin{enumerate}
    \item Si $kp < n$~: $W^{k,p}(\Omega) \hookrightarrow\hookrightarrow L^q(\Omega)$
          pour tout $q < p^*$ (injection compacte).
    \item Si $kp > n$~: $W^{k,p}(\Omega) \hookrightarrow\hookrightarrow C(\overline{\Omega})$
          (injection compacte).
\end{enumerate}
\end{theoreme}

\begin{exemple}[Cas $n = 1$]
En dimension $1$, pour $p = 2$ et $k = 1$~: $kp = 2 > n = 1$, donc
$H^1(\Omega) \hookrightarrow C(\overline{\Omega})$. Toute fonction $H^1$ en
dimension $1$ est continue (apr\`es modification sur un ensemble de mesure nulle).
\end{exemple}

\begin{exemple}[Cas $n = 3$, $k = 1$, $p = 2$]
$kp = 2 < n = 3$, donc $p^* = \frac{3 \cdot 2}{3 - 2} = 6$.
On a $H^1(\Omega) \hookrightarrow L^6(\Omega)$ (continu) et
$H^1(\Omega) \hookrightarrow\hookrightarrow L^q(\Omega)$ pour $q < 6$ (compact).
\end{exemple}

\section{Th\'eor\`eme de trace}

\begin{theoreme}[Trace]
\label{thm:trace}
\index{th\'eor\`eme de trace}
Soit $\Omega$ un ouvert born\'e \`a bord $C^1$. L'op\'erateur de restriction
$\gamma_0 : C^\infty(\overline{\Omega}) \to C^\infty(\partial\Omega)$,
$\gamma_0(u) = u|_{\partial\Omega}$, se prolonge en un op\'erateur lin\'eaire
continu surjectif~:
\[
    \gamma_0 : W^{1,p}(\Omega) \to W^{1-1/p, p}(\partial\Omega),
    \quad 1 < p < \infty.
\]
De plus, $\ker \gamma_0 = W_0^{1,p}(\Omega)$.
\end{theoreme}

\begin{remarque}
Ce th\'eor\`eme justifie la condition de Dirichlet $u = g$ sur $\partial\Omega$
dans le cadre des espaces de Sobolev~: on demande $\gamma_0(u) = g$ au sens de
la trace, m\^eme si $u$ n'est pas continue.
\end{remarque}

\section{In\'egalit\'e de Poincar\'e}

\begin{theoreme}[In\'egalit\'e de Poincar\'e]
\label{thm:poincare}
\index{in\'egalit\'e de Poincar\'e}
Soit $\Omega \subset \R^n$ un ouvert born\'e et $1 \leq p < \infty$. Il existe
une constante $C = C(\Omega, p, n) > 0$ telle que~:
\[
    \norm{u}_{L^p(\Omega)} \leq C \norm{\nabla u}_{L^p(\Omega)}
    \quad \forall\, u \in W_0^{1,p}(\Omega).
\]
\end{theoreme}

\begin{proof}
Par l'absurde. Supposons qu'il existe une suite $(u_k)$ dans $W_0^{1,p}(\Omega)$
avec $\norm{u_k}_{L^p} = 1$ et $\norm{\nabla u_k}_{L^p} \to 0$. Par Rellich-Kondrachov,
$(u_k)$ admet une sous-suite convergente dans $L^p$. La limite $u$ satisfait
$\nabla u = 0$ p.p., donc $u$ est constante. Comme $u \in W_0^{1,p}$, on a $u = 0$,
ce qui contredit $\norm{u}_{L^p} = 1$.
\end{proof}

\begin{corollaire}
Sur $W_0^{1,p}(\Omega)$ avec $\Omega$ born\'e, la semi-norme
$\norm{\nabla u}_{L^p}$ est \'equivalente \`a la norme $\norm{u}_{W^{1,p}}$.
\end{corollaire}

\begin{theoreme}[In\'egalit\'e de Poincar\'e-Wirtinger]
\index{in\'egalit\'e de Poincar\'e-Wirtinger}
Soit $\Omega$ un ouvert born\'e connexe \`a bord lipschitzien. Pour
$1 \leq p < \infty$~:
\[
    \norm{u - \bar{u}}_{L^p(\Omega)} \leq C \norm{\nabla u}_{L^p(\Omega)}
    \quad \forall\, u \in W^{1,p}(\Omega),
\]
o\`u $\bar{u} = \frac{1}{\abs{\Omega}} \int_\Omega u\, dx$ est la moyenne de $u$.
\end{theoreme}

\begin{formules}[title=\textbf{R\'esum\'e des injections de Sobolev ($k=1$)}]
\begin{center}
\renewcommand{\arraystretch}{1.3}
\begin{tabular}{lcl}
\textbf{Condition} & \textbf{Injection} & \textbf{Type} \\
\hline
$p < n$ & $W^{1,p} \hookrightarrow L^{np/(n-p)}$ & continue \\
$p < n$ & $W^{1,p} \hookrightarrow\hookrightarrow L^{q}$, $q < \frac{np}{n-p}$ & compacte \\
$p = n$ & $W^{1,p} \hookrightarrow L^q$, $\forall q < \infty$ & continue \\
$p > n$ & $W^{1,p} \hookrightarrow C^{0,1-n/p}$ & continue \\
$p > n$ & $W^{1,p} \hookrightarrow\hookrightarrow C(\overline{\Omega})$ & compacte \\
\end{tabular}
\end{center}
\end{formules}

\section{Exercices}

\begin{exercice}
Montrer que $u(x) = \abs{x}^{-\alpha}$ appartient \`a $W^{1,p}(B(0,1))$
(boule unit\'e de $\R^n$) si et seulement si $(\alpha + 1)p < n$.
\end{exercice}

\begin{exercice}
Soit $\Omega = (0,1)$ et $u(x) = x^\beta$ pour $\beta > 0$. D\'eterminer pour
quelles valeurs de $\beta$ et $k$ on a $u \in H^k(\Omega)$.
\end{exercice}

\begin{exercice}
D\'emontrer l'in\'egalit\'e de Sobolev en dimension $1$~: pour $u \in W^{1,1}(\R)$,
on a $\norm{u}_{L^\infty} \leq \frac{1}{2}\norm{u'}_{L^1}$.
\end{exercice}

\begin{exercice}
Montrer que $H^1(\R^n)$ ne s'injecte pas dans $L^\infty(\R^n)$ pour $n \geq 2$.
\emph{Indication~:} construire un contre-exemple logarithmique.
\end{exercice}

\begin{exercice}
Soit $\Omega$ un ouvert born\'e de $\R^n$ \`a bord lisse. Montrer que la constante
de Poincar\'e $C_P$ v\'erifie $C_P \leq \frac{\operatorname{diam}(\Omega)}{2}$ quand $p = 2$.
\end{exercice}

\begin{exercice}
Montrer que si $u \in H^1_0(\Omega)$ et $f \in L^2(\Omega)$ v\'erifient $-\Delta u = f$
au sens faible, alors $\norm{u}_{H^1_0} \leq C\norm{f}_{L^2}$ avec une constante
ne d\'ependant que de $\Omega$.
\end{exercice}
